\documentclass[10pt,a4paper]{article}

% --- Packages ---
\usepackage[utf8]{inputenc}
\usepackage[T1]{fontenc}
\usepackage{lmodern}
\usepackage[margin=0.8in]{geometry}
\usepackage{microtype}
\usepackage{booktabs}
\usepackage{tabularx}
\usepackage{multirow}
\usepackage{xcolor}
\usepackage{listings}
\usepackage{enumitem}
\usepackage{tcolorbox}
\usepackage{fancyhdr}
\usepackage[hidelinks,colorlinks=true,linkcolor=navyblue,citecolor=navyblue,urlcolor=navyblue]{hyperref}

% --- Colors ---
\definecolor{navyblue}{RGB}{16, 44, 87}
\definecolor{slate}{RGB}{53, 64, 86}
\definecolor{lightgrey}{RGB}{245, 247, 250}
\definecolor{bordergrey}{RGB}{210, 215, 225}
\definecolor{codebg}{RGB}{240, 243, 246}
\definecolor{accentbox}{RGB}{235, 242, 250}

% --- Page Setup ---
\pagestyle{fancy}
\fancyhf{}
\lhead{\small\textcolor{slate}{\textbf{Illumo} --- Architecture Consensus (Unified)}}
\rhead{\small\textcolor{slate}{2026-08-04}}
\cfoot{\thepage}
\renewcommand{\headrulewidth}{0.4pt}

\setlist[itemize]{noitemsep, topsep=3pt, leftmargin=1.5em}
\setlist[enumerate]{noitemsep, topsep=3pt, leftmargin=1.5em}

% --- Code Listings ---
\lstset{
    backgroundcolor=\color{codebg},
    basicstyle=\ttfamily\small,
    breaklines=true,
    captionpos=b,
    commentstyle=\color{gray},
    keywordstyle=\color{navyblue}\bfseries,
    stringstyle=\color{teal},
    frame=single,
    rulecolor=\color{bordergrey},
    framerule=0.5pt,
    xleftmargin=3pt,
    xrightmargin=3pt
}

\begin{document}

% --- Title Block ---
\begin{center}
    {\LARGE \bfseries \textcolor{navyblue}{Illumo --- Architecture Consensus (Unified)}} \\[0.4em]
    {\small \textbf{Status:} Single living document --- \textbf{authoritative for later sessions} \quad|\quad \textbf{Last updated:} 2026-08-04}
\end{center}

\vspace{0.2em}
\noindent This file \textbf{merges and supersedes} scattered design memory into one coherent story. Read this first; treat external PDFs and old agenda notes as \textbf{history} (\S 2).

\vspace{0.8em}

\begin{table}[h!]
\centering
\small
\begin{tabularx}{\textwidth}{l X}
\toprule
\textbf{Source} & \textbf{Role in this document} \\
\midrule
Desktop ``Main agenda for Illumo'' & Product wishlist $\rightarrow$ \S7 scorecard \\
\texttt{Illumo\_Architecture\_Decisions.pdf} (ChatGPT) & CA / OOP-practice design review $\rightarrow$ purpose, ownership, input, rules, SYCL \\
\texttt{Illumo\_Engine\_Architecture\_Decisions.pdf} (ChatGPT) & Aspirational 2D engine $\rightarrow$ what we kept vs discarded \\
\texttt{gpt\_illumo\_arch\_assessment.pdf} (ChatGPT) & Assessment of \textit{current} tree $\rightarrow$ risks, direction \\
Grok architecture reviews + local code review & Strengths, bugs, debt \\
In-repo LaTeX design notes + decision log & Formal decision IDs (D-*) \\
\texttt{docs/current\_issues.txt} & Product/correctness punch list \\
\textbf{Current code under \texttt{Illumo/Source/}} & \textbf{Wins} when anything conflicts \\
\bottomrule
\end{tabularx}
\end{table}

\begin{tcolorbox}[colback=lightgrey,colframe=slate,boxrule=0.5pt,arc=2pt]
\textbf{Rule:} If this document and the code disagree, \textbf{code wins} until this file is updated in the same change set.
\end{tcolorbox}

\noindent Optional deeper reading (not required to resume work):
\begin{itemize}
    \item \texttt{docs/illumo-design-notes.tex} / \texttt{.pdf} --- long-form chapters
    \item \texttt{docs/sections/09-design-decision-log.tex} --- append-only formal decision prose
\end{itemize}

\vspace{0.5em}
\tableofcontents
\newpage

% --- Section 0 ---
\section{One-line summary}
\textbf{Illumo is a cellular-automata learning sandbox in a modular-monolith shell: App owns composition, Illumo owns services, modules drive sim/UI, rendering is enroll-once + token stream (OpenGL + Mock), Canvas is a dense grid with RGB fade display. Keep that shape. Fix Start/load/Wireworld correctness and docs drift. Defer ECS, multi-pass pipelines, multi-backend, infinite chunks, and SYCL until product pain or learning goals demand them.}

% --- Section 1 ---
\section{Project purpose and scope}

\subsection{Primary purpose (still true)}
From the CA design review and the original spirit of the repo:
\begin{itemize}
    \item \textbf{Practice OOP and systems programming}, not ship a commercial game engine.
    \item The main learning sandbox is \textbf{cellular automata} (GoL family, multi-state rules, editor, camera, console).
    \item Custom allocators, Tracy, backend interfaces, etc. are \textbf{useful experiments}, not requirements that must dominate the design.
    \item Do \textbf{not} rewrite everything for an imagined perfect architecture. Refactor where coupling produces \textbf{real friction}.
\end{itemize}

\subsection{What the project grew into}
It is no longer ``one Game of Life file.'' It is a \textbf{small CA framework / app}:
\begin{itemize}
    \item Multiple rule sets, edit/run modes, save/load hooks
    \item Token-based OpenGL presentation + headless tests
    \item Engine-shaped packaging (App / Engine / Game / Rendering / Services)
\end{itemize}
That growth is \textbf{acceptable}. Architecture should reflect \textbf{actual families of behavior} (CA + UI + host), not a fantasy multi-genre engine.

\subsection{Explicit non-goals (for now)}
\begin{itemize}
    \item Production AAA renderer or full multiplayer engine
    \item Full ECS / archetypes / cached structural queries
    \item Scene/World graph as the primary cell path
    \item Second real graphics API (Vulkan/Metal) as a near-term deliverable
    \item Perfect Linux/macOS parity before Windows remains solid
    \item Infinite canvas + SYCL \textbf{before} dense-grid product paths are correct and documented
    \item Aggressive batching/instancing/render graphs without profiling evidence
\end{itemize}

\subsection{Allocators (from CA design PDF --- retained policy)}
Retain arena / stack / pool allocators as learning code and optional utilities. Do \textbf{not} force them into every service. Appropriate uses:

\begin{table}[h!]
\centering
\small
\begin{tabularx}{\textwidth}{l X}
\toprule
\textbf{Allocator} & \textbf{Fit} \\
\midrule
Arena & Frame scratch, transient command data, parsers, load buffers \\
Stack & Nested temporary lifetimes (LIFO free) \\
Pool & Many same-sized objects with stable slots \\
\bottomrule
\end{tabularx}
\end{table}

A general-purpose allocator (mimalloc-class) is a different product; do not reinvent it.

% --- Section 2 ---
\section{Historical sources (how to read old notes)}
Three generations of ``design truth'' got mixed in chat history. This section untangles them so later sessions do not re-implement superseded plans.

\subsection{Desktop agenda (``Main agenda for Illumo'')}
Early product wishlist:
\begin{table}[h!]
\centering
\small
\begin{tabularx}{\textwidth}{l X}
\toprule
\textbf{Item} & \textbf{Intent} \\
\midrule
On-screen text / fonts & UI \\
In-game command line & Tools \\
Console drawn in the GL window & UI \\
More rule sets & Content \\
Infinite canvas in \textbf{16$\times$16 chunks} + hash map & Scale \\
Mouse drag pan & UX \\
SYCL (or similar) for large cell counts & Scale / learning \\
\bottomrule
\end{tabularx}
\end{table}

\noindent \textbf{Outcome today:} UI/tools/pan/most rulesets largely done. \textbf{Chunk infinite canvas + SYCL not done} (deliberately deferred). Full scorecard $\rightarrow$ \S 7.

\subsection{CA design review (\texttt{Illumo\_Architecture\_Decisions.pdf})}
Themes that \textbf{still guide} the project:
\begin{table}[h!]
\centering
\small
\begin{tabularx}{\textwidth}{l X}
\toprule
\textbf{Theme} & \textbf{Guidance} \\
\midrule
Ownership & Explicit Engine/App object; reduce uncontrolled global lookup \\
Input & Thin layer: GLFW callbacks $\rightarrow$ InputManager $\rightarrow$ application logic \\
Boundaries & Contain GLFW/OpenGL behind wrappers; no GL in sim/UI domain \\
Sim vs view & Separate simulation from presentation; double-buffer generations \\
Rules & Families (life-like / elementary / multi-state); JSON for \textbf{data}, not executable behavior \\
Parallel & Parallelize \textbf{simulation data} last---after serial works \\
SYCL & Optional compute backend; never a hard engine dependency \\
Refactor plan & Ownership $\rightarrow$ input extract $\rightarrow$ controllers $\rightarrow$ contain GL $\rightarrow$ rule families $\rightarrow$ double-buffer $\rightarrow$ then parallel \\
\bottomrule
\end{tabularx}
\end{table}

\noindent \textbf{Principle:} \textit{Abstract real volatility; hard-code stable structure; ownership explicit.} \quad (Closer to current Illumo than the ``engine passes'' PDF.)

\subsection{Engine design PDF (\texttt{Illumo\_Engine\_Architecture\_Decisions.pdf})}
Aspirational \textbf{general 2D engine} design (July 2026). Valuable as boundary thinking; \textbf{not} the current product plan.

\begin{table}[h!]
\centering
\small
\begin{tabularx}{\textwidth}{l X}
\toprule
\textbf{Proposed} & \textbf{Status in Illumo} \\
\midrule
Illumo owns services; context as view & \textbf{Kept} \\
Backend owns GPU resources + ID tables & \textbf{Kept} (IBackend / GL registries) \\
Modules as subsystems & \textbf{Kept} \\
RenderPipeline with Opaque / Transparent / UI / Debug \textbf{passes} & \textbf{Not built} \\
MeshDrawCommand-style typed pass queues & \textbf{Not built} --- we use tagged-union \texttt{RenderCommand} stream \\
Minimal EntityTable (transform, mesh, texture) & \textbf{Archived} (D-E3) \\
Transform hierarchies + fixed-tick mesh interpolation & \textbf{Not needed} for CA; cell \textbf{fade} is the visual interpolation \\
Full ECS (sparse-set sketches) & \textbf{Deferred forever-until-pain} \\
Multi-backend, render graphs, heavy batching & \textbf{Explicitly deferred} in that PDF too \\
\bottomrule
\end{tabularx}
\end{table}

\noindent \textbf{Do not} re-introduce pass objects or entity mesh tables ``because the engine PDF said so.'' Token stream + drawable list is the shipped boundary.

\subsection{Later assessments (Grok + ChatGPT on the \textit{current} tree)}
Agreement across reviews:
\begin{itemize}
    \item Layered modular monolith is \textbf{appropriate}
    \item App owns modules; token path + MockBackend are real strengths
    \item Scene-as-list is correct; dead graph/EntityTable should stay gone
    \item Main risks: incomplete module Start handling, Canvas dual role, partial backend neutrality, and command-queue overflow
    \item Highest value: product correctness + one canonical Canvas model in docs --- \textbf{not} ECS or multi-pass
\end{itemize}

% --- Section 3 ---
\section{Overall assessment}
Illumo is a \textbf{coherent layered modular monolith}:

\begin{lstlisting}[language=bash]
Platform
  -> App/CellMain          (product composition)
  -> Engine/Illumo         (owns services + module list)
  -> Modules               (CellGameModule, DebugModule)
  -> Game + Rulesets       (CA domain / rules)
  -> Rendering             (Scene list, drawables, Renderer, tokens)
  -> IBackend              (GLBackend | MockBackend)
\end{lstlisting}

It is an \textbf{engine-shaped application}, not a cleanly separated ``engine product + game product.'' That is \textbf{intentional at current scale}.

\noindent \textbf{Verdict:} Functional and reasonably testable. Residual risk is mostly \textbf{product correctness holes} and \textbf{documented debt}, not a need for ECS or a full rewrite.

% --- Section 4 ---
\section{Strengths (keep these)}
\begin{table}[h!]
\centering
\small
\begin{tabularx}{\textwidth}{l X}
\toprule
\textbf{Area} & \textbf{Consensus} \\
\midrule
\textbf{App owns composition} & \texttt{CellMain}: \texttt{Init} $\rightarrow$ \texttt{addModule} $\rightarrow$ \texttt{StartModules} (D-E1). Engine does not construct Game modules. \\
\textbf{Module lifecycle} & \texttt{Start} / \texttt{Update} / \texttt{DispatchDrawables} / \texttt{Exit} is clear. \\
\textbf{Render split} & Enroll once; emit tokens per frame; backend executes (D-R1--D-R8, D-R10). \\
\textbf{Rulesets} & Strategy hierarchy; pure \texttt{nextState} + \texttt{evalCell}; double-buffered generation (D-P3). \\
\textbf{Scene model} & Per-frame drawable list only (D-E3, D-E4). \\
\textbf{Tests} & Headless \texttt{IllumoTests}: rules, canvas, tokens, UI, renderer inject. \\
\textbf{Debt hygiene} & Dead experiments under \texttt{archive/} rather than half-live. \\
\bottomrule
\end{tabularx}
\end{table}

% --- Section 5 ---
\section{Canonical architecture (code truth)}

\subsection{Packages}
\begin{table}[h!]
\centering
\small
\begin{tabularx}{\textwidth}{l X}
\toprule
\textbf{Package} & \textbf{Role} \\
\midrule
\textbf{App/} & Product composition (\texttt{CellMain}). \\
\textbf{Engine/} & Host: Illumo, modules, frozen \texttt{IllumoContext}. \\
\textbf{Game/} & CA domain + presentation (\texttt{Canvas}, \texttt{CellGameModule}, \texttt{CellContext}). \\
\textbf{Rulesets/} & CA rules (GoL family, Wireworld, \dots). \\
\textbf{Rendering/} & Scene list, drawables, Renderer, tokens, OpenGL, Mock. \\
\textbf{Services/} & Log, env, input, console UI, save/load API, allocators. \\
\textbf{Foundation/} & Macros, math aliases (\texttt{MathTypes.h}), sysinfo. \\
\textbf{Platform/} & OS entry + native dialogs. \\
\textbf{Assets/} & Asset loaders. \\
\textbf{Tests/} & Headless suite. \\
\bottomrule
\end{tabularx}
\end{table}

\noindent \textbf{House style (D-008 / CONTRIBUTING):} avoid \texttt{auto}; avoid namespaces (prefer static classes/structs); no recursion; third-party via PR --- unless a later decision waives.

\subsection{Ownership and lifetime}
\begin{itemize}
    \item \textbf{Illumo} owns long-lived services with \texttt{unique\_ptr} (window, renderer, camera, env, input, scene, command line, \dots).
    \item \textbf{IllumoContext} is a \textbf{non-owning} pointer bag for modules (D-E5: frozen for the two shipped modules; validate at \texttt{Start}).
    \item \textbf{App} decides which modules exist.
    \item Do not grow the bag for convenience; a third module with different needs should take \textbf{explicit constructor dependencies}.
    \item Prefer explicit references for domain logic over ``look up everything from the bag.''
\end{itemize}

\subsection{Module contract}
\begin{lstlisting}[language=C++]
class IModule {
public:
    virtual void Start(IllumoContext* context) = 0;
    virtual void Update(double dt) = 0;
    virtual void DispatchDrawables(Scene* scene) = 0;
    virtual void Exit() = 0;
};
\end{lstlisting}

\begin{table}[h!]
\centering
\small
\begin{tabularx}{\textwidth}{l X}
\toprule
\textbf{Hook} & \textbf{Responsibility} \\
\midrule
Start & Bind inputs, allocate game state; bail if context incomplete \\
Update & Simulation / input \\
DispatchDrawables & Contribute what should draw this frame \\
Exit & Teardown module-owned state \\
\bottomrule
\end{tabularx}
\end{table}

\noindent Shipped modules: \texttt{CellGameModule} (always), \texttt{DebugModule} (\texttt{\#ifndef NDEBUG}).\\
\textbf{Known hole:} \texttt{Start} can fail without preventing later \texttt{Update} / \texttt{DispatchDrawables} $\rightarrow$ \S 8.

\subsection{Frame loop (production)}
\begin{lstlisting}[language=bash]
CellMain
  -> Illumo::Init()                    // services only
  -> addModule(CellGameModule)
  -> addModule(DebugModule)            // Debug builds
  -> StartModules()
  loop:
    Update(dt)   // InputManager -> Camera -> modules.Update
    Render()
      scene.ClearDrawables()
      modules.DispatchDrawables(scene)
      renderer.BeginFrame()
      renderer.RenderScene(scene, camera)   // tokens, then optional immediate fallback
      renderer.EndFrame()                   // swap (GL)
\end{lstlisting}

\noindent Typical draw order from the game module: canvas $\rightarrow$ console $\rightarrow$ FPS $\rightarrow$ splash.

\subsection{Rendering architecture (shipped)}
\textbf{Mental model:} Game never issues draw \texttt{gl*} for the live path. It enrolls resources and emits \textbf{tokens}.

\begin{lstlisting}[language=bash]
Game / Rulesets / UI
    |  (no gl* for draw submission)
    v
Drawable::AppendCommands(Renderer*)
    |  enroll handles + push tokens
    v
CommandQueue<RenderCommand>     // tagged union payloads
    |
    v
IBackend::SubmitCommandQueue
    |
    +-- GLBackend + GLDevice   (handle resolve, PBO upload, bind tracker)
    +-- MockBackend            (tests)
\end{lstlisting}

\begin{table}[h!]
\centering
\small
\begin{tabularx}{\textwidth}{l X}
\toprule
\textbf{Piece} & \textbf{Job} \\
\midrule
\textbf{Modules} & Choose what should appear this frame. \\
\textbf{Scene} & Ordered list of drawable pointers for this frame only. \\
\textbf{Drawable} & Prefer \texttt{AppendCommands(Renderer*)} $\rightarrow$ tokens. Immediate \texttt{Draw()} only if AppendCommands returns false (tests/stubs). \\
\textbf{Renderer} & Frame setup tokens; collect drawable tokens; submit; then hybrid immediate list. \\
\textbf{IBackend} & Create resources, queue, submit, begin/end frame. \\
\textbf{GLBackend / GLDevice} & Real OpenGL: handle registries, execute tokens, PBO texture updates, bind-state tracking, blend-func-on-enable (D-R5). \\
\textbf{MockBackend} & Headless: record creates + command order. \\
\bottomrule
\end{tabularx}
\end{table}

\noindent \textbf{Enroll (rare):} \texttt{enrollMesh} / \texttt{enrollShader} / \texttt{enrollTexture} $\rightarrow$ opaque \texttt{unsigned long} table IDs.\\
\textbf{Per frame:} \texttt{RenderCommand} tagged union (bind, uniform, update texture/buffer, draw, clear, viewport, pipeline).\\
\textbf{Not current (old engine PDF):} separate Opaque/Transparent/UI pass \textit{objects}, MeshDrawCommand-only world draws, entity mesh tables.\\
\textbf{Production pure-token drawables (D-R10):} Canvas, CommandLine, GLString, SplashText.\\
\textbf{Token migration phases 0--6:} complete for planned scope.

\subsection{Canvas presentation model (canonical --- code wins)}
\textbf{Important history:} D-P4 briefly moved presentation to \textbf{GPU R8 state texture + 256$\times$1 palette} (snap colors, drop CPU fade). Fade was later \textbf{restored}. Intermediate LaTeX notes that still describe ``R8-only current'' are \textbf{stale}.

\noindent \textbf{Code truth after fade restore:}
\begin{table}[h!]
\centering
\small
\begin{tabularx}{\textwidth}{l X}
\toprule
\textbf{Layer} & \textbf{Contents} \\
\midrule
\textbf{Domain} & \texttt{lifeCanvas} --- dense \texttt{unsigned char[width $\times$ height]} (\textbf{not} chunked) \\
\textbf{View} & CPU palette $\rightarrow$ \texttt{targetRgb} / \texttt{displayRgb} fade; \texttt{texCanvasBuffer} RGB bytes \\
\textbf{GPU} & RGB display texture; dirty-rect \texttt{UpdateTexture} (PBO); shader samples \texttt{uDisplayTexture} \\
\bottomrule
\end{tabularx}
\end{table}

\noindent \textbf{What survived from D-P4:} dirty-rect uploads, PBO path in \texttt{GLTexture::UpdateSubImage}, bind-state tracker, dirty flags for idle frames (D-P1).\\
\textbf{What was rolled back:} exclusive R8+palette GPU path as the live presentation; dual float RGB staging + fade are back.\\
\textbf{Decision D-C1:} Keep Canvas as intentional domain + view + GPU enroll \textbf{monolith} until pure-sim tests or large grids force \texttt{LifeGrid} / \texttt{CanvasView}.

\noindent \textbf{Scale walls:}
\begin{itemize}
    \item \texttt{calcGeneration} always walks the \textbf{full} grid (rect args ignored).
    \item Neighbor counts: scalar Moore/toroidal.
    \item Change detection may do a second full pass for dirty AABB (helps sparse \textbf{upload}, still O($W \times H$) sim).
    \item Per cell: 1 B life + $\sim$24 B float RGB + 3 B display tex. Default 80$\times$60 is fine; large/infinite needs redesign.
\end{itemize}

\subsection{Rules and encoding}
\textbf{Active rules} (factory / AllSets): Game of Life, Seeds, Brian's Brain, Highlife, Day \& Night, Life Without Death, \textbf{Wireworld}.\\
\textbf{Stubs:} Rule 90 / 184 (identity \texttt{nextState}).

\begin{lstlisting}[language=C++]
class RuleSet {
public:
    virtual unsigned char nextState(unsigned char cell,
                                    unsigned char aliveNeighbors) const = 0;
    virtual void evalCell(const unsigned char& target,
                          unsigned char dest[3]) const = 0; // palette colors
    void calcGeneration(...); // double-buffer: nextGen then write-back
};
\end{lstlisting}

\begin{table}[h!]
\centering
\small
\begin{tabularx}{\textwidth}{l X}
\toprule
\textbf{Encoding} & \textbf{Values} \\
\midrule
Binary / life-like & \texttt{0} = alive, \texttt{1} = dead; neighbor count treats value \texttt{0} as live \\
Multi-state (e.g. Brian's Brain) & $\ge 2$ additional states (e.g. dying = 2) \\
\textbf{Wireworld} & \texttt{0} head, \texttt{1} empty, \texttt{2} tail, \texttt{3} conductor (head = 0 reuses head-neighbor counting) \\
\bottomrule
\end{tabularx}
\end{table}

\noindent \textbf{Generation path (D-P3):}
\begin{enumerate}
    \item Count neighbors from raw \texttt{lifeCanvas} (toroidal).
    \item Write \texttt{nextState} into scratch \texttt{nextGen}.
    \item If any cell changed: copy back, mark sparse dirty AABB for visuals. Still life $\rightarrow$ no dirty.
\end{enumerate}

\noindent Rules stay free of rendering and input. \texttt{evalCell} supplies palette/RGB colors only.

\noindent \textbf{Optional cleanup (from CA PDF):} collapse life-like rules into one family + JSON birth/survive tables:
\begin{lstlisting}
{ "family": "life_like", "birth": [3], "survive": [2, 3] }
\end{lstlisting}
JSON describes \textbf{data}, not executable behavior. Multi-state (Wireworld, Brian's Brain) stay separate families.

\subsection{Simulation vs rendering}
\begin{enumerate}
    \item Simulation advances on a \textbf{tps $\times$ speedFactor} clock (accumulator + step cap).
    \item Rendering runs every frame.
    \item Double-buffer prevents reading a half-written generation.
    \item Visual fade is \textbf{presentation}, not simulation state.
    \item Modes in \texttt{CellGameModule}: \texttt{NORMAL | EDIT | EXIT}; save/load are registered commands, not frame states.
\end{enumerate}

\subsection{Input}
Intended flow: \texttt{GLFW callbacks / poll $\rightarrow$ InputManager $\rightarrow$ module / controller logic}.\\
\textbf{Today:} InputManager holds key/mouse state and contexts; CellGameModule / DebugModule consume it.\\
\textbf{D-E2:} InputManager must not depend on Game types.

\subsection{Developer console}
The Debug-only console keeps general tooling in \symbolref{CommandLine} while
\symbolref{CellGameModule} registers simulation, canvas, camera, ruleset, and
save/load behavior through \symbolref{CommandRegistry}. Registry usage,
description, and completion metadata drive live help and Tab completion.
Registered commands no longer fall through as unknown. Dense saves are validated
on load, activate their saved ruleset, copy row-by-row into the current canvas,
and rebuild the visible state immediately. The old \texttt{vid\_restart} no-op is
not advertised because safe GL-context recreation requires full resource
re-enrollment.

\subsection{Window / platform boundaries}
\begin{itemize}
    \item Semantic window ops hide GLFW types from the main loop.
    \item OpenGL calls stay under \texttt{Rendering/OpenGL/} (+ window bootstrap).
    \item Interfaces only where multiple implementations or volatility are real (\texttt{IBackend}, \texttt{IRenderWindow}).
    \item macOS Metal myths: GLFW can create a no-client-API window and expose native handles.
\end{itemize}

% --- Section 6 ---
\section{Locked decisions (catalog)}
Full formal prose lives in \texttt{docs/sections/09-design-decision-log.tex}.

\subsection{Structure and style}
\begin{table}[h!]
\centering
\small
\begin{tabularx}{\textwidth}{l X}
\toprule
\textbf{ID} & \textbf{Decision} \\
\midrule
\textbf{D-001} & Package rename: Core$\rightarrow$Game, System$\rightarrow$\{App,Engine,Services\}, Init+Util$\rightarrow$Foundation, AssetLoaders$\rightarrow$Assets. \\
\textbf{D-002} & Math aliases in \texttt{MathTypes.h} (not \texttt{Math.h} --- Windows CRT collision). \\
\textbf{D-004} & Modules: Start / Update / DispatchDrawables / Exit. \\
\textbf{D-005} & Each CA variant is a \texttt{RuleSet} subclass (factory in CellContext; registry later optional). \\
\textbf{D-006} & Modes = \texttt{CellState} enum in module (not archived State class hierarchy). \\
\textbf{D-008} & House style: avoid \texttt{auto}; no namespaces; no recursion; third-party via PR. \\
\textbf{D-N1} & Adopt Illumo as the project, product, CMake-target, runtime, and save-format name; keep the existing \texttt{Illumo} host class. \\
\bottomrule
\end{tabularx}
\end{table}

\subsection{Rendering (D-R*)}
\begin{table}[h!]
\centering
\small
\begin{tabularx}{\textwidth}{l X}
\toprule
\textbf{ID} & \textbf{Decision} \\
\midrule
\textbf{D-003 / D-R*} & Token submission via IBackend + RenderCommand; game must not permanently depend on raw GL for draw. \\
\textbf{D-R1} & \texttt{RenderCommand} = simple tagged union (C-style), not byte-stream compiler / \texttt{std::variant}. \\
\textbf{D-R2} & Each drawable emits tokens via \texttt{AppendCommands}; Renderer owns frame setup + submit. \\
\textbf{D-R3} & Opaque \texttt{unsigned long} handles + backend registries for v1; strong typedefs deferred. \\
\textbf{D-R4} & Migration phases 1--6 done; destroy resources at shutdown only for v1; Windows GL primary. \\
\textbf{D-R5} & On blend enable, always set \texttt{glBlendFunc} (console panel transparency). \\
\textbf{D-R6} & Archive dead render queue / SceneObject graph helpers; Scene + Renderer is the path. \\
\textbf{D-R7} & MockBackend + headless tests; no Vulkan/Metal yet. \\
\textbf{D-R8} & Inject \texttt{IBackend*} into Renderer for tests; production owns GLBackend. \\
\textbf{D-R9} & String-named uniforms = debt if a second \textit{real} GPU API appears. \\
\textbf{D-R10} & Production drawables pure-token; hybrid \texttt{Draw()} for tests/stubs only. \\
\textbf{D-007} & Enroll resources outside the per-frame stream (frame queue = bind/draw/update). \\
\bottomrule
\end{tabularx}
\end{table}

\subsection{Performance (D-P*)}
\begin{table}[h!]
\centering
\small
\begin{tabularx}{\textwidth}{l X l}
\toprule
\textbf{ID} & \textbf{Decision} & \textbf{Note} \\
\midrule
\textbf{D-P1} & Dirty visual path --- idle frames skip full recolor/upload. & Still current \\
\textbf{D-P2} & UI batch: CommandLine one packed update; GLString geometry cache. & Still current \\
\textbf{D-P3} & Double-buffer \texttt{calcGeneration} + sparse dirty AABB. & Still current \\
\textbf{D-P4} & Originally: R8 + palette + dirty-rect PBO; drop dual float RGB. & \textbf{Partially superseded:} RGB fade display restored \\
\bottomrule
\end{tabularx}
\end{table}

\subsection{Engine shape (D-E*, D-C*, D-F*)}
\begin{table}[h!]
\centering
\small
\begin{tabularx}{\textwidth}{l X}
\toprule
\textbf{ID} & \textbf{Decision} \\
\midrule
\textbf{D-E1} & Module registration lives in App; Engine knows only \texttt{IModule}. \\
\textbf{D-E2} & InputManager has no Game types. \\
\textbf{D-E3} & EntityTable archived; cells are not entities. \\
\textbf{D-E4} & Scene is drawable list only (no graph). \\
\textbf{D-E5} & Freeze IllumoContext; validate at Start; third module $\rightarrow$ explicit deps. \\
\textbf{D-C1} & Canvas dual role intentional until scale forces split. \\
\textbf{D-F1} & MacroDefs / Windows.h include toxicity deferred until real pain. \\
\bottomrule
\end{tabularx}
\end{table}

% --- Section 7 ---
\section{Early agenda vs today}
\begin{table}[h!]
\centering
\small
\begin{tabularx}{\textwidth}{l X}
\toprule
\textbf{Agenda item} & \textbf{Status} \\
\midrule
Render text / fonts & \textbf{Done} (FreeType + GLString / SplashText) \\
Command line & \textbf{Done} (commands can still grow) \\
Console on GL screen & \textbf{Done} (token UI) \\
More rulesets & \textbf{Partly} --- Wireworld live; 90/184 stubs \\
Infinite 16$\times$16 chunk canvas & \textbf{Not done} --- dense grid by design for now \\
Mouse pan & \textbf{Done} (camera controls) \\
SYCL / large-grid parallel & \textbf{Not done} --- serial full-grid; deferred \\
\bottomrule
\end{tabularx}
\end{table}

% --- Section 8 ---
\section{Known product / correctness issues}

\subsection{Bugs (high signal)}
\begin{table}[h!]
\centering
\small
\begin{tabularx}{\textwidth}{c l X}
\toprule
\textbf{\#} & \textbf{Severity} & \textbf{Issue} \\
\midrule
1 & bug & \textbf{Partial Start still crashes later.} If \texttt{IllumoContextHasGameCore} fails, \texttt{Start} returns without \texttt{cellContext}, but \texttt{Update} / \texttt{DispatchDrawables} still use it $\rightarrow$ null deref. \\
2 & bug & \textbf{Wireworld editor can't place heads with the mouse.} \texttt{setcell} can place numeric states, but there is no convenient head/tail brush. \\
3 & bug & \textbf{Startup seed is always a GoL glider (value 0).} Under WIREWORLD that plants five electron heads, not a wire. \\
\bottomrule
\end{tabularx}
\end{table}

\subsection{Suggestions}
\begin{table}[h!]
\centering
\small
\begin{tabularx}{\textwidth}{c X}
\toprule
\textbf{\#} & \textbf{Issue} \\
\midrule
4 & Wireworld tests miss 2-head birth / 3-head no-birth; save/load command integration coverage remains limited. \\
\bottomrule
\end{tabularx}
\end{table}

\subsection{Assessment-only risks (structural)}
\begin{itemize}
    \item Command queue \textbf{2048} capacity with \textbf{silent drop}; raw pointer payloads must stay alive until submit.
    \item CMake source/config duplication was resolved on 2026-08-02 with shared source lists and a common target-configuration helper; empty package CMake files remain.
    \item Renderer.h still includes GL types for production construction.
    \item Hybrid token + immediate Draw path remains for stubs.
\end{itemize}

\subsection{Looks fine (do not ``fix'' as bugs)}
\begin{itemize}
    \item Wireworld encoding (head = 0) and head-neighbor counting
    \item Fade order for paint/sim path
    \item \texttt{rebuildPalette} $\rightarrow$ full dirty refresh on ruleset switch
    \item SceneObject removal has no leftover live call sites
    \item Pure-token production drawables
\end{itemize}

% --- Section 9 ---
\section{Architectural debt (do when it hurts)}
\begin{table}[h!]
\centering
\small
\begin{tabularx}{\textwidth}{l X}
\toprule
\textbf{Item} & \textbf{Notes} \\
\midrule
String uniforms & D-R9 --- GL-shaped until second real backend \\
Hybrid Draw path & Tests/stubs only; production is tokens \\
Command queue silent overflow & Prefer log/assert if near cap \\
CMake duplication & Resolved 2026-08-02 with shared source lists/settings; no forced package libraries \\
Renderer.h includes GL types & Factory / pimpl later \\
Full-grid sim + float fade memory & Scale wall for large canvases \\
MacroDefs + Windows.h & D-F1 deferred \\
IllumoContext growth & Frozen; third module = explicit deps \\
Life-like JSON family collapse & Optional cleanup of repetitive RuleSet classes \\
Chunk + SYCL & Optional after serial benchmark / product need \\
Resource destroy/reload & Shutdown-only for v1 (D-R4); hot-reload later \\
\bottomrule
\end{tabularx}
\end{table}

% --- Section 10 ---
\section{Recommended work order}

\subsection{A. Correctness first}
\begin{enumerate}
    \item Failed \texttt{Start} $\rightarrow$ do not run that module's Update/Dispatch.
    \item Wireworld seed + mouse head-placement UX.
    \item Add focused save/load command integration tests.
    \item Keep this file and LaTeX ``current state'' sections aligned with \textbf{RGB fade Canvas}.
\end{enumerate}

\subsection{B. Hygiene}
\begin{enumerate}[resume]
    \item CMake source/config consolidation --- \textbf{completed 2026-08-02}.
    \item Command-queue overflow policy (log/assert).
\end{enumerate}

\subsection{C. Only if product or learning goals require it}
\begin{enumerate}[resume]
    \item Canvas split (\texttt{LifeGrid} / \texttt{CanvasView}).
    \item Chunked / infinite canvas (agenda 16$\times$16 map).
    \item Threaded or SYCL simulation backend behind a narrow interface.
    \item Backend factory + non-string uniforms.
    \item Data-driven life-like rule family (JSON birth/survive).
    \item Focused controllers if input coupling becomes painful.
\end{enumerate}

\noindent \textbf{Explicitly deferred:} Full Scene/World abstractions, general ECS, render graphs, multi-backend, aggressive batching.

% --- Section 11 ---
\section{Core design principles (merged)}
\begin{enumerate}
    \item \textbf{CA learning sandbox first} --- not a general engine product.
    \item \textbf{Ownership explicit} --- Illumo owns services; App owns module set; context does not own.
    \item \textbf{Boundaries over cleverness} --- abstract volatility; hard-code stable structure.
    \item \textbf{Sim produces complete state; render observes} --- double-buffer; no draw mid-generation.
    \item \textbf{Tokens for draw submission} --- enroll once, emit commands, backend executes.
    \item \textbf{Parallelize data transformations last} --- serial correctness first; chunks/SYCL optional.
    \item \textbf{Archive experiments} --- don't leave half-live ECS/graph/passes in the hot path.
    \item \textbf{Simplest architecture that preserves the boundaries you care about}.
    \item \textbf{Code wins over docs} --- update this file when consensus shifts.
\end{enumerate}

% --- Section 12 ---
\section{Open questions (remaining)}
\begin{table}[h!]
\centering
\small
\begin{tabularx}{\textwidth}{l X}
\toprule
\textbf{Topic} & \textbf{Working answer} \\
\midrule
Resource ownership long-term & Destroy only at shutdown for v1 (D-R4). Refcounts / generational handles later if hot-reload needs them. \\
Linux/macOS parity & Freeze until Windows token path solid --- Windows path is solid; parity still optional. \\
Tracy CI policy & Debug-oriented; no strict CI policy yet. \\
When to introduce chunks / SYCL & After correctness + serial benchmarks, or as an explicit learning goal. \\
\bottomrule
\end{tabularx}
\end{table}

% --- Section 13 ---
\section{How to use this document}
\begin{table}[h!]
\centering
\small
\begin{tabularx}{\textwidth}{l X}
\toprule
\textbf{When} & \textbf{Action} \\
\midrule
New chat / new session & Read \textbf{this file first}. \\
Deep dive & Then LaTeX design notes / decision log. \\
Architecture change & Update \textbf{this file} + append a decision log entry in \texttt{09-design-decision-log.tex}. \\
Code lands & Update ``code truth'' sections here if behavior changed. \\
Old PDFs / agenda & Treat as \textbf{history} (\S 2); do not re-implement superseded engine plans by default. \\
Bug triage & \S 8 first; architecture is not the problem until proven otherwise. \\
\bottomrule
\end{tabularx}
\end{table}

\newpage
\appendix
\section{Appendix A --- Source map (where code lives)}
\begin{table}[h!]
\centering
\small
\begin{tabularx}{\textwidth}{l X}
\toprule
\textbf{Concern} & \textbf{Typical location} \\
\midrule
Frame loop / composition & \texttt{Source/App/CellMain.cpp} \\
Host / services / modules & \texttt{Source/Engine/Illumo.*} \\
CA module / modes & \texttt{Source/Game/CellGameModule.*}, \texttt{CellContext.*} \\
Grid + fade + GPU enroll & \texttt{Source/Game/Canvas.*} \\
Rules & \texttt{Source/Rulesets/*} \\
Tokens / Renderer & \texttt{Source/Rendering/Renderer.*}, \texttt{RenderCommand.*} \\
GL execute & \texttt{Source/Rendering/OpenGL/*} \\
Mock & \texttt{Source/Rendering/Mock/*} \\
Console & \texttt{Source/Services/CommandLine.*} \\
Dead experiments & \texttt{archive/dead-engine/}, \texttt{archive/dead-render/}, \texttt{archive/old-states/} \\
\bottomrule
\end{tabularx}
\end{table}

\section{Appendix B --- What not to do next}
\begin{itemize}
    \item Do not reintroduce EntityTable / SceneObject graph ``for completeness.''
    \item Do not build Opaque/Transparent/UI pass objects unless profiling or a real multi-genre product appears.
    \item Do not make SYCL a hard dependency of Illumo or Game.
    \item Do not document R8-only presentation as current without verifying shaders + Canvas.cpp.
    \item Do not expand IllumoContext casually for a third module.
\end{itemize}

\end{document}
